

\subsection{The spherical pairs}
Let $F$ be a non-Archimedean local field of characteristic zero.
We first consider the linear pair
\begin{equation}\label{eq:main-linear-subgroup}
 G_{\mathrm L}=\GL_4(F)\times\GL_2(F),\qquad
 H_{\mathrm L}=\{(\diag(a,b),b):a,b\in\GL_2(F)\}.
\end{equation}
\begin{mainthm}[Linear spherical pair]\label{thm:main-linear}\label{m1:thm:main}
For every irreducible admissible representation $\pi$ of $G_{\mathrm L}$,
\[
 \dim\Hom_{H_{\mathrm L}}(\pi,\mathbb C)\leq1.
\]
\end{mainthm}

For the unitary statement let $E/F$ be a quadratic field extension,
write $z\mapsto\bar z$ for its nontrivial automorphism, and put
$E^1=\{z\in E^\times:z\bar z=1\}$. Forms are linear in the first
variable. For a Hermitian space $Z$, write $\U(Z)$ and $\GU(Z)$
for its groups of $F$-rational isometries and similitudes.
Let $V_1,V_2$ be hyperbolic Hermitian planes. Write $\lambda$ for
every similitude multiplier. With $V=V_1\perp V_2$, define
\begin{equation}\label{eq:main-unitary-subgroup}
 \begin{split}
 G_{\mathrm U}&=\GU(V)\times\GU(V_2),\\
 H_{\mathrm U}&=\{(u\oplus h,h):u\in\GU(V_1),\ h\in\GU(V_2),
                              \ \lambda(u)=\lambda(h)\}.
 \end{split}
\end{equation}
Since $V$ is hyperbolic,
$G_{\mathrm U}\simeq\GU(2,2)(F)\times\GU(1,1)(F)$.
In the antidiagonal realization the two planes are the outer and middle
coordinate planes; the fixed permutation in
Section~\ref{m2:red:pair} identifies that realization with
\eqref{eq:main-unitary-subgroup}.
\begin{mainthm}[Unitary similitude spherical pair]\label{thm:main-unitary}\label{m2:thm:main}
For every irreducible admissible representation $\pi$ of $G_{\mathrm U}$,
\[
 \dim\Hom_{H_{\mathrm U}}(\pi,\mathbb C)\leq1.
\]
\end{mainthm}
These are Models~1 and~2 among the strongly tempered spherical pairs
studied by Wan--Zhang~\cite{WanZhang25}. They computed the unramified local relative characters and established multiplicity formulas in the tempered setting. For the linear pair, global questions have also been
studied by Pollack--Wan--Zydor~\cite{PWZ21} and by
Cauchi--Gutierrez Terradillos~\cite{CGTLinear}, including connections
with central $L$-values.

\subsection{Unitary similitude Shalika uniqueness}
Fix a nontrivial additive character $\psi:F\to\mathbb C^\times$.
Let $W=X\oplus Y$ be the split skew-Hermitian four-space with
$X=Y=E^2$ and matrix
$J_W=\left(\begin{smallmatrix}0&I_2\\-I_2&0\end{smallmatrix}\right)$.
Let $B$ be a hyperbolic Hermitian form on $X$, with matrix
$\mathsf B=\mathsf B^\dagger$, where $a^\dagger=\bar a^{\,t}$.
Write $\lambda_B$ for the multiplier on $\GU(B)$ and $\tr_E$ for
the matrix trace over $E$. The embedding $m_B:\GU(B)\to\GU(W)$,
the Shalika subgroup, and its character are
\begin{equation}\label{eq:intro-shalika}
 \begin{gathered}
 N(X)=\left\{n(u)=\begin{pmatrix}I_2&u\\0&I_2\end{pmatrix}:
                                    u^\dagger=u\right\},\\
 m_B(a)=\diag\bigl(a,\lambda_B(a)\bar a^{-t}\bigr),\qquad
 S_B^{\GU}=m_B(\GU(B))N(X),\\
 \psi_B(n(u))=\psi\bigl(-\tfrac12\tr_E(u\mathsf B)\bigr).
 \end{gathered}
\end{equation}
The trace $\tr_E(u\mathsf B)$ belongs to $F$. Representations of
$\GU(B)$ are viewed as representations of $m_B(\GU(B))$ through
this embedding.
\begin{mainthm}[Unitary similitude Shalika models]\label{thm:main-gu}\label{gu:thm:main}
For every irreducible admissible representation $\Pi$ of $\GU(W)$
and every irreducible admissible representation $\sigma$ of $\GU(B)$,
\begin{equation}\label{gu:eq:main}
 \dim\Hom_{S_B^{\GU}}(\Pi,\sigma\boxtimes\psi_B)\leq1.
\end{equation}
\end{mainthm}

For ordinary unitary groups, Gan~\cite{GanPeriods} and Chen--Gan
\cite[Proposition~3.4 and Theorem~3.5]{ChenGan} establish the theta-period
identity for Shalika periods with arbitrary coefficient representations. The distinction between
isometry and similitude groups matters: Shalika periods for ordinary unitary groups can have multiplicity greater than one
\cite[Theorem~4.3]{ChenGan}. Here we establish the corresponding period identity
for unitary similitudes and use it to prove Theorem~\ref{thm:main-gu}.

The twisted Jacquet module of the similitude oscillator gives a
Hom-space identity (Theorem~\ref{gu:thm:hom}) relating the
Shalika model to the unitary spherical pair. Combined with
Theorem~\ref{thm:main-unitary}, it bounds the Shalika Hom space of a
big theta lift by one. To pass to an arbitrary irreducible admissible
$\GU(W)$-representation $\Pi$, Theorem~\ref{gu:asm:coverage} provides a
unitary character $\chi$ of $E^\times/F^\times$ and
$\tau\in\Irr(\GU(V))$ with
\begin{equation}\label{gu:eq:intro-coverage}
 \Theta^{\GU}_{\chi,1,1}(\tau)\twoheadrightarrow\Pi.
\end{equation}
The surjection in \eqref{gu:eq:intro-coverage} induces,
by precomposition, an injection from the Shalika Hom space
of $\Pi$ into that of the big theta lift.
Theorem~\ref{thm:main-gu} follows from the
multiplicity-one bound for the latter.
The coverage theorem uses analytic equal-rank theta
dichotomy and highly ramified stability to obtain
nonvanishing in the prescribed quasi-split tower.

\subsection{Outline of the proofs}
Section~\ref{sec:quiver} develops the determinant-open quiver geometry
shared by the two pairs: closed orbits, stabilizers, and normal
representations. Section~\ref{sec:linear} establishes the linear
two-leg theorem and treats the identity descendant.
These results, together with generalized
Harish--Chandra descent, yield Theorem~\ref{thm:main-linear}.

For the unitary pair, this common geometry is only the starting point.
Section~\ref{sec:unitary-geometry} determines the Hermitian domains
admitting rational points, proves uniqueness of the rational orbit
inside each admitted closed geometric orbit, and deduces rational
transpose stability. It then determines the corresponding \(F\)-rational normal representations and their corrected odd operators. Section~\ref{sec:unitary-descendants}
shows that stabilizer-invariant distributions on these
normal spaces are fixed by the corrected
involutions. The split-adjoint regular case
reduces to the linear two-leg theorem, while
the anisotropic case is treated using Hermitian
Gram fibers. The scalar case requires a separate
identity-descendant argument. This proves
Theorem~\ref{thm:main-unitary}.

Section~\ref{sec:gl-shalika} recovers Rallis's multiplicity-one
theorem for twisted Jacquet modules of $\GL_4(F)$, as recorded by
Prasad~\cite[Theorem~1]{PrasadShalika}, using the linear spherical-pair
theorem and Gan's theta-period identity~\cite[Theorem~3.1]{GanPeriods}.
Section~\ref{sec:gu-transfer} establishes the corresponding Shalika
period identity for unitary similitude theta lifts. Finally,
Section~\ref{sec:gu-coverage} shows that, after varying the splitting
character, every irreducible admissible representation of $\GU(W)$ is
a quotient of a suitable big theta lift. This completes the proof of
Theorem~\ref{thm:main-gu}.

Appendices~A and~B supply the functional-analytic and normalization
details for the theta-period arguments. Appendix~C treats general
splitting data, while Appendix~D records the rational transpose
transporter constructions used in the unitary orbit analysis.

\subsection{Notation}
Unless otherwise stated, representations of $p$-adic groups
are smooth and defined over $\mathbb C$.
We write $\pi^\vee$ for the smooth contragredient of $\pi$
and $\omega_\pi$ for its central character when $\pi$
is irreducible admissible.
We denote by $\Irr(G)$ the set of isomorphism classes
of irreducible admissible representations of $G$. The symbols $G_{\mathrm L},H_{\mathrm L}$ and
$G_{\mathrm U},H_{\mathrm U}$ denote the two spherical pairs.
In the theta part $G_\theta=\GU(V)$ and
$G'_\theta=\GU(W)$ denote the two members of the dual pair;
$K=\GU(B)$ denotes the Shalika coefficient group.
All similitudes scale their forms by $\lambda$ in the direction
$h(gv,gw)=\lambda(g)h(v,w)$. We distinguish ordinary transpose $t$, conjugate
transpose $\dagger$, and the form adjoint $*$.
Absolute values satisfy $|z|_E=|N_{E/F}z|_F$ and
$|s|_E=|s|_F^2$ for $s\in F^\times$.
Schwartz spaces consist of compactly supported locally constant
functions, and tensor products are algebraic. Ordinary and compact induction are written
$\Ind$ and $\cInd$; their normalization and duality are specified
where used.

\subsection*{Acknowledgments}
The author thanks Professor Dihua Jiang for helpful suggestions on an earlier
draft of this paper.

\subsection*{Disclosure of AI Use}
ChatGPT (including GPT-5.6 Sol and GPT-6 Astra) and Claude Opus 5.5 were used for proof exploration and drafting this
paper. The author takes responsibility for the
mathematical content.